\documentclass{ieeeaccess}
\usepackage{cite}
\usepackage{amsmath,amssymb,amsfonts}
\usepackage{graphicx}
\usepackage{textcomp}
\usepackage{booktabs}
\usepackage{multirow}
\usepackage{bm}

\def\BibTeX{{\rm B\kern-.05em{\sc i\kern-.025em b}\kern-.08em
    T\kern-.1667em\lower.7ex\hbox{E}\kern-.125emX}}

\begin{document}
\history{Date of publication xxxx 00, 0000, date of current version xxxx 00, 0000.}
\doi{10.1109/ACCESS.2024.0429000}

\title{Unlocking the Potential of Pseudo-Labels: A Visual Anchoring Framework for Cross-Domain Behavior Annotation}
\author{\uppercase{First Author}, \uppercase{Second Author}}
\address[1]{Institution}
\corresp{Corresponding author: First Author (e-mail: author@institution.edu).}

\begin{abstract}
As behavioral analysis across application domains scales up, the demand for annotated video data grows exponentially, presenting a significant bottleneck. This study introduces an automated visual anchoring framework utilizing general-purpose Vision-Language Models (VLMs) and Large Language Models (LLMs) to reliably annotate out-of-distribution video datasets. Specifically, we investigate whether generic human behavior representations can correctly guide annotation on specialized domains. By breaking down annotation tasks into single-frame conceptual alignment and multi-frame temporal reasoning, we evaluate zero-shot annotation generation without domain-specific re-training. 
Our findings confirm that when visual behaviors are sufficiently grounded in distinct frame-level states, pseudo-labels reach near-ground-truth quality. Pseudo-label fine-tuned models achieved 88.65\% accuracy on BowTurnHead and 43.22\% on TeacherBehavior. Selective annotation routing further mitigates distribution shifts, improving robust behavior recognition.
\end{abstract}

\begin{keywords}
Action Recognition, Large Language Models, Pseudo-Labeling, Video Understanding
\end{keywords}

\titlepgskip=-21pt
\maketitle
\section{Introduction}
\PARstart{T}{his} paper demonstrates the utility of VLMs in cross-domain annotation. 
% Please wait while I migrate the rest of the text safely.
\EOD
\end{document}
